\section{Limits of the model}
\label{sec:discussion}

This section takes the limits of the model itself, in the order in which they could change the comparison between the rigs. \Cref{sec:disc_details} gives the sources and numbers behind several of the estimates.

\subsection{Canopy life and kite handling}
\label{sec:disc_kite}

Neither the life of a ship kite's canopy nor how reliably it can be launched and recovered has been measured. Both bear most on the kite's place against the deck rigs.

\textbf{Canopy life.} No ship kite has a measured life. No maker states more than \num{4000} flight hours for any kite or sail, and the derived range of \num{1200} to \num{6900} cannot be narrowed from what is published (\cref{sec:canopy_life}). Developers of small kites expected much less, neither measured: about \num{1000} hours for a \SI{25}{m^2} leading-edge inflatable kite \citep{vandervlugt2013}, and not much more than 100 hours at full load for light ram-air wings of 15 to \SI{21}{m^2} \citep{bormann2013}. We keep the low bound of \num{1200} hours because those cloths are far lighter and held at full load, which a ship kite rarely is. SkySails, the German maker whose kites flew on several cargo ships, calls fatigue a limit on economic success and gives no number \citep{paulig2013}. Within that range the life reverses kite and rotors on all three ships studied (\cref{sec:res_missed}), which no other input does at its bounds in \cref{tab:inputs}. A measured canopy life of a ship kite would narrow the comparison more than any other number.

\textbf{Launch and recovery.} The model charges each launch and landing for crew time, retries, tears, losses and turns of the ship into a suitable heading, and gives no force while it lasts. On the \SI{140}{m} ship a cycle costs \tbd{USD 235 to 286} at the goal for tears and losses. \tbd{All launch inputs at their high bound together take 1.2 to 6.7 points of cost saving there and put the best kite behind the rotors on every route.} The model assumes that launches are possible day and night whenever wind and sea allow. Experience makes this its weakest assumption. Launching by day only takes \tbd{0.2 to 4.0 points, median 0.7}, since a kite up at dusk stays up through the night. SkySails limits operation to non-freezing conditions \citep{paulig2013}. We do not model this. It could matter only on winter legs at high latitude, and we have not measured how much.

The record of ship kites comes from SkySails (\cref{sec:disc_details_kites}). Launch and landing worked, but the supplier reported limited launch opportunities in heavy sea in its first test phase, and components that held up poorly at sea. Operators said the system could only be run in daylight and was too labour-intensive for rotating crews, the crews that change every few months \citep{cedelft2016}. The founder later said the handling system would have to be optimised \citep{wrage2018}.

\subsection{Rig reliability}
\label{sec:disc_reliability}

A rig that is out of service saves nothing but still costs its capital. No rig has an independently measured availability, failure rate or repair time, so we set one goal of 3\,\% out of service for all three. It comes from the one long statement by an operator rather than a maker. The rotors of the ferry Estraden ``could be used 97\% of the time'' over about six years \citep{spliethoff2021}, without saying of what total time. We read it as a share per unit and per round trip while the wind would allow use, which the source does not say. Makers report more: above 99\,\% on the same ship \citep{kuuskoski2018}, above 95\,\% on the bulk carrier Afros over eight years \citep{anemoi2026}, and 99.6\,\% for the wings of the cargo ship Canop\'ee over two years \citep{oceanwings2025}. All of these are undefined and on small rigs or first ships. For a ship kite there is no figure at all, and the 3\,\% is our assumption, so that a missing record neither credits nor penalises it. The upper bound of 15\,\% is one unsourced estimate of ours for all three rigs, two outages of 20 days in 250 sea days. Low use of the first kites was a mix of wind, handling, economics and owners stopping \citep{cedelft2016}, which is not the same as being out of service.

Class rules, the technical rules of the classification societies that certify a ship, treat the rig as auxiliary. A ship with a rig down still makes its route on its engines, and the model routes it so. Stopped rotors stand as drag. The coefficient of 0.64 from our fit at zero spin lies above the 0.5 usually assumed \citep{rechevilanova2021}, with no full-scale measurement of either, and a hazard study warns that a rotor stopped for about a day risks its bearing \citep{emsa2023}. The model has units fail independently, while a blackout or a shared controller would stop them all, and a rig whose units fail together falls behind the next at a lower share (\cref{sec:res_missed}). Repair is left in the rig maintenance fraction of 2\,\%, applied alike to all rigs. One study assumes 5\,\% for kites and 2\,\% for other rigs \citep{emsa2023}. \tbd{Raising the kite from 3 to 5\,\% took 0.09 to 0.38 points from it without changing the order, measured when every rig stood at 3\,\%. The final run repeats this from 2\,\%.} We add no cost for schedule or for a kite recovered from the water, as no record gives one. As our estimate, one extra port call a year for a repair would cost any rig 0.2 to 0.6 points.

\subsection{Series prices}
\label{sec:disc_prices}

All three rigs are priced at one production stage, a series of 50 or more units (\cref{sec:costs}): manufacturing cost from the masses and costs per kilogram of comparable industries, times a market factor of 2, the ratio of selling price to manufacturing cost. Only wings and rotors are made in series today. For rotors this lands within 13\,\% below and 6\,\% above Norsepower's 2019 prices \citep{tillig2020b}, and the makers of wings and rotors have built dozens of units. For the kite the gap is wider. SkySails sold about ten ship kites \citep{tonline2012}, as whole systems at about EUR~0.5~million for \SI{160}{m^2} in 2007 \citep{wiwo2008} and about EUR~1~million for \SI{320}{m^2} \citep{wintecc2009_record,tonline2012}. That is \tbd{2.6 to 7} times our price for the same systems\tbd{, compared at the sizes as sold; at their projected areas our price is 10 to 15\,\% lower, so these ratios rise by about that}. Unit cost falls with the number built, and the learning rate is that fall per doubling. Onshore wind power has become cheaper by 12\,\% on average per doubling of installed capacity, with single studies from $-11$ to 32\,\% \citep{rubin2015}. From about ten kites sold, an optimistic 15 to 20\,\% per doubling reaches our price only after \tbd{200 to \num{40000}} units, and the wind-power mean after \tbd{\num{1800} to \num{380000}}. Our kite price is a goal of series production that no kite has reached. A check against kite power is in \cref{sec:disc_details_prices}.

\subsection{Rigs larger than any built}
\label{sec:disc_size}

On the largest ships the cost optimum lies at rigs larger than any built: rows of wings and rotors that fill the deck, and kites far larger than any that has flown. The model then extrapolates beyond its evidence.

\textbf{Deck rigs.} No ship in service carries more than five rotors \citep{gcaptain2021seazhoushan,maritimeexecutive2024sohar}. The first with six, a ro-ro, is fitting out \citep{windtech2026toulouse}. The wings and rotors are modelled at the same level as the kite: steady force coefficients, interference between neighbours, and the same leeway, heel and rudder balance. Interference, the loss of force of a unit in the disturbed wind of its neighbours, is the least certain of these coefficients and matters most for arrays this long. Flow simulations and wind-tunnel tests of two to four units find losses of up to 6 to 37\,\% of the driving force, and gains of up to 17\,\% downwind (\cref{sec:disc_details_size}). The rotor tests stop at a spin ratio of 2, below the 3 at which our rotors run, while the deflection of the wind behind a rotor grows with the spin ratio. Our array factor, the share of the force of isolated units that a row keeps, depends on the number and spacing of the units, averaged over the wind angle. Against CFD (computational fluid dynamics) of wing sails on a bulk carrier it lies low for widely spaced sails and high for close ones. Both rigs take deck area that a cargo ship may need for cargo or cranes, which the model does not price. Both rigs can stay up when the kite would have to come down, which the model counts only as the absence of a launch limit. On these ships the comparison is between the largest deck rig that fits and a kite that needs no deck.

\textbf{Where the rigs stand.} A rig's side force turns the ship unless it acts where the hull's own side force acts, and the rudder holds any moment left over at the cost of drag. A single wing or rotor stands at the hull's lateral centre, about a third of the length forward of midship, where its side force needs no rudder. A row spreads over the deck and is centred at midship, so the rudder carries the yaw moment. How much this matters depends on which foil, hull or rudder, makes side force at less drag, and the literature differs. On a cargo ship modelled from CFD the hull was about three times as efficient as the rudder, and moving the sails' centre of effort from midship to its best position changed the resistance by only 1 to 2\,\% \citep{kramer2022}. ShipCLEAN finds the rudder the better foil and gains several points of fuel saving by moving rotors aft \citep{tillig2020}. We use the effective rudder aspect ratio of about 1.6 that \citet{kramer2022} derive from CFD, rather than the doubling that is often assumed because the hull above the rudder acts as a mirror plane.

\textbf{Large kites.} The cost optimum lies at kite areas that have never flown, and on the Megamax, the Neo-Panamax and the VLCC at the \SI{16000}{m^2} upper bound of the search, where a larger kite might save more of the cost (\cref{sec:res_size}). A \SI{3200}{m^2} kite of aspect ratio 3 spans about \SI{98}{m}. Two numbers decide whether it can exist: the wing loading it can take and the mass of its canopy. We hold the canopy at the cloth mass of the method (\cref{sec:kite}) and the wing loading at \SI{0.5}{kN/m^2} of projected area, and scale everything else. SkySails designed kites of 225 and \SI{320}{m^2} for 160 and \SI{320}{kN} \citep{paulig2013}, 0.7 and \SI{1.0}{kN/m^2}, and reports wing loads of up to \SI{1}{kN/m^2} from operation \citep{skypower100_tib}. These are about 0.9 to \SI{1.3}{kN/m^2} projected if their areas are flat. We keep the lower value as a choice, because the canopy's life and the rig's strength are stated at it and no source gives how long a canopy lasts at a higher one. A kite built for a higher limit saves less of the cost on each of three ships, and one built for \SI{0.3}{kN/m^2} about a point more (\cref{sec:res_force}), so \SI{0.5}{kN/m^2} is not the limit most favourable to the kite. The largest ram-air wings built are cargo parachutes of \SI{900}{m^2} and \SI{1040}{m^2}, the heaviest at \SI{600}{kg}, sewn in modules \citep{dunker2013}. They fly far below our wing loading and once at a time. The largest ship kite is \SI{320}{m^2} flat \citep{fritz2013}. A wider span is harder to launch and to handle, which the model does not count. A large wing keeps its area per line by adding rows of lines, up to eight \citep{dunker2013}, while our bridle keeps its count and thickens its lines, so its drag is independent of size. The turn radius rests on kites up to \SI{320}{m^2} flat and may grow with the mass per area \citep{erhard2013b}. \Cref{sec:res_aspect} gives what a higher aspect ratio would be worth if it could be handled at the same cost.

\textbf{VTOL kites.} The result for the VTOL kite rests on one wing that was designed and never built, MX2 of the company Makani, scaled over more than a factor of ten in area. Its goals are far from the record: a life of 15 years against the 22~h that the M600, Makani's \SI{600}{kW} prototype, flew. Evidence pulls in both directions. The M600 came out 40 to 70\,\% over its mass targets and showed far more parasitic drag, the drag of all parts that make no lift, than MX2 assumes. Its one flight from a floating base ended with the loss of the kite in hover (\cref{sec:rigid_kite}). So the base may be too kind. Against that, its prices are those of a series of 50 or more, and the wing keeps the aspect ratio of MX2, designed to generate and not to pull. The favourable case moves mass, drag, propeller area and prices to their bounds together, but not the aspect ratio, so a wing designed for towing could gain more than we show. We tested two ships and not all thirteen classes, so a larger or faster ship could come out differently. The gap in price per unit of thrust does not depend on the ship. The wing's inertia around the loop is neglected, tested against flight only on Makani's \SI{60}{kg} Wing 7, where the lowest wind comes out about \SI{1}{m/s} high (\cref{sec:rigid_kite}). Several kites share one loop at a momentum loss, from flying through air already slowed by the others. Little or no flight backs either of these two parts. The shared loop bears only on fleets, so not on the base result, in which one small wing comes out best.

\textbf{The kite keeps its design shape in a quasi-steady eight.} Quasi-steady means that the forces at each point of the figure eight are those of steady flight at that point, with the canopy keeping its designed shape. We add a drag coefficient for the losses of a real canopy, calibrated on a \SI{25}{m^2} kite. The glide ratio of 3.85 with the tether lies between what earlier ship-kite studies assume, 3.5 \citep{traut2014} to 4.6 \citep{podeur2016}, and at the 3.9 fitted to the \SI{160}{m^2} SkySails kite \citep{erhard2013b}. Lift coefficients above 0.8 are hard to reach with textile kites \citep{penloup2021}. Whether a \SI{3200}{m^2} kite loses the same, more or less is unknown. \tbd{Without its flight drag the kite saves 1.4 to 2.0 points more of the cost on e-ammonia and 0.6 to 0.8 points more on VLSFO, on kites of \SIlist{800;3200}{m^2} on the \SIlist{140;200}{m} ships; from an earlier run.}

\subsection{Where the kite tows from}
\label{sec:tow_point}

The kite's side force acts where its tether meets the ship. Ahead of the hull's lateral centre it turns the bow to leeward, aft of it the stern, and the rudder must hold the difference at the cost of its drag (\cref{sec:balance}). We tow every soft kite from its launch mast at the forward end of the deck, ahead of the lateral centre by 2 to 10\,\% of the length on every class. To see what that position costs, we flew each class's cheapest kite on e-ammonia with its tow point at the bow, at the lateral centre and at midship, on the westerly and trade-wind routes (\cref{fig:tow_point}). Only the rudder load changes between them in the model. The mast, its price and the cargo it displaces stay the same.

\begin{figure*}[tbp]
\centering
\plot[\textwidth]{tow_point}
\caption{Fuel (a) and cost (b) saving of each class's kite with its tow point at the hull's lateral centre, at midship with the mast taking the deck slots it takes at the bow, the least it can take there, and at midship with the mast taking one whole deck bay of a container ship, the most, less the saving with its tow point at the bow, the line at zero; the marker is the mean over the two routes and the bar their range. Classes are sorted by deadweight.}
\label{fig:tow_point}
\end{figure*}

The bow costs almost no thrust. The lateral centre saves \tbd{0.3} points more fuel than the bow on average, less than the \tbd{$-1.6$ to $+1.5$} points by which the two differ between ships and routes as the router takes a different path. Midship saves \tbd{1.2} points less fuel than the bow on average, as its side force acts about a third of the length aft of the lateral centre. The loss is \tbd{0.5} points on the bulk carriers and tankers and \tbd{1.9} on the other classes. Fuel is only part of the cost, so the cost saving moves by \tbd{0.1} points on average and at most \tbd{1.3}.

Cargo space decides the rest. On a bulk carrier or tanker a mast at midship stands between the hatches or beside the cargo piping and takes no cargo space, but it gains nothing in thrust either. On a container ship it stands in the deck stack. No container ship flies a kite, so we estimate the loss from the geometry: at most one 40-foot bay across the beam on every tier on deck, and at least the single row of a bay that the bow mast takes (\cref{sec:model}). With the bay and the rudder load together, midship saves \tbd{1.9 to 4.7} points less cost than the bow on the container ships that fill their slots before their deadweight. The bow is therefore never behind midship by more than the routing's spread. We keep it, as the kites built so far do \citep{fritz2013}, and give up the \tbd{0.3} points of fuel saving of the lateral centre, which lies within the cargo area on most ships.

\subsection{One wind profile}
\label{sec:disc_profile}

The wind grows with height, and the kite flies hundreds of metres higher than the deck rigs stand, so the assumed growth decides how much more wind the kite gets. All rigs take their wind from one power-law profile with exponent 0.05, which gives 1.18 times the \SI{10}{m} wind at \SI{270}{m}, where the base kite flies. It matches the mean of ERA5 model-level profiles over the study region \citep{hersbach2020}, but the shear, the growth of the wind with height, differs by region (\cref{sec:disc_details_profile}). It is weaker in the trade-wind belt and stronger north of \SI{40}{\degree N}. One profile therefore understates the kite's thrust on the North Atlantic routes and overstates it in the trades, by 10 to 20\,\%. The kite's cost saving on the route to New York is therefore more likely understated than overstated. Other ship studies use exponents of 0.11 \citep{rechevilanova2021}, 1/7 \citep{penloup2021} and 0.27 \citep{tillig2020}, above the mean shear that ERA5 or the measurements show over the open sea, and so give the kite more wind than we do.

\subsection{Limits that act on every rig}
\label{sec:disc_common}

The remaining limits act on the ship with and without a rig, and on all rigs alike, so they move the cost savings more than the order of the rigs.

\textbf{Six routes, three years of weather.} The routes cover the trade winds of both hemispheres, the westerlies, the Bay of Biscay and the African coast, but all start in Rotterdam and lie in the Atlantic. They stand for wind climates, not for the trades of each class. A Capesize does not call at Bilbao. The Pacific and the Indian Ocean are missing, and three years cannot show the variation between years. The wind field is coarse, \SI{1.5}{\degree} and six hours. Weather at \SI{1}{\degree} or six hours, or coarser, makes the routed voyage look worse than it is, most in winter and most for ships that rely on the wind \citep{wendling2025}, so our wind field more likely understates the gain from wind than overstates it.

\textbf{Calm-water resistance of full hulls.} The calm-water resistance is the base against which every fuel and cost saving is measured. Each class hull is calibrated so that its calm-water power at service speed matches its source design (\cref{sec:model}). The resistance comes from the method of \citet{holtrop1982}, a regression on model tests. The Ultramax and Kamsarmax bulk carriers, with prismatic coefficients of 0.87 and 0.89, lie beyond the hulls it covers. Evaluated there unchanged it would make their calm-water power at operating speed too low, by 7 and 18\,\%, so we evaluate the formulas at 0.85 for these two hulls, with their true displacement (\cref{sec:disc_details_hull}). No measured speed-power curve of such a hull checks this choice.

\textbf{Added resistance in waves is a mean estimate.} The recalibrated formula we use explains 73\,\% of the variance of \num{1127} model-test samples from 25 ships, against 67\,\% before recalibration. Its errors are largest in short waves and near resonance, where the wave length is close to the ship's length \citep{mittendorf2022}. It acts alike on the ship with and without a rig, except where heavy seas ground the kite.

\textbf{The wind on the ship itself.} Measured wind loads of single ships scatter about the wind-tunnel coefficients of their type \citep{blendermann1993}, and the bulk carriers borrow the laden tanker's shape without their cranes and hatch coamings, the raised rims of the cargo hatches. The method matters most where the deck cargo is large. Against a simpler estimate with no side force, a drag on the beam times a \SI{10}{m} deckhouse at $C_D=0.7$, it lowers the cost saving of each rig on fuel oil by 0.06 to 0.34 points on the four container ships and the car carrier. The order of the rigs that save money stays the same.

\textbf{Round trips on fixed trades.} Every ship sails out and back between the same two ports, with the return passage free to bend into the westerlies or the trades within its corridor. That is how liner services run. Tramp bulk carriers and tankers, which have no fixed route, instead follow the next cargo, often on triangular or ballast voyages, and a wind-assisted tramp operator could prefer voyages that follow the clockwise circulation of the North Atlantic winds. Our round trips understate what such an operator could gain.
